% Exact research-candidate source-aware proof for P06.1.
% Not a live-book input before independent and production verification.
% No novelty or priority claim is made.

\begin{proof}[Proof of \Cref{lem:p061-nilpotent-rpl-invariance}]
The assertion is local for the G-topology.  Work first on a strict affinoid
chart
\[
                  X=\mathcal M(A),\qquad
                  X_0=\mathcal M(A/I),
\]
where \(I^e=0\) for some \(e\).  Every multiplicative seminorm on \(A\)
kills \(I\), so restriction gives a bijection
\[
                         \mathcal M(A/I)\longrightarrow\mathcal M(A).
\]
It is a homeomorphism: the source and target are compact Hausdorff, the map
is continuous, and it is bijective.  Equivalently, the topology can be
read directly from the evaluation functions, whose values are unchanged
modulo \(I\).

Rational domains correspond under this homeomorphism.  Indeed, lift all
functions occurring in a rational-domain presentation from \(A/I\) to
\(A\).  Their inequalities cut out the same set because seminorms kill
\(I\).  If the reduced functions generate the unit ideal, their lifts do
as well: an identity modulo \(I\) gives an element \(1+n\), with
\(n\in I\), in the lifted ideal, and \(1+n\) is a unit by a finite
geometric series.  Rational localization commutes with the quotient and
again has nilpotent kernel.  By the Gerritzen--Grauert theorem in the form
recalled in \cite[\S11.1.6]{CLD25}, every affinoid domain is a finite union
of rational domains.  Hence this identifies the rational bases and
therefore the induced G-topologies of \(X_0\) and \(X\).  The same argument
on affinoid charts proves the asserted global identification.

At the sheaf level, \cite[Lemma~5.4(3)]{GJR26} states exactly that
\(\operatorname{PL}_{\Lambda}\) is invariant under passage to the
reduction, for every coefficient group \(\Lambda\subset\mathbf R\).  The
reductions of \(X_0\) and \(X\) are canonically the same because the
defining ideal of \(X_0\) is nilpotent.  Applying that lemma to both spaces
already proves the asserted equivalence for \(\Lambda=\mathbf R\).  For
completeness, the direct unit-lifting calculation is as follows.  Pullback
preserves \(\mathbf R\)-PL functions by
\cite[Lemma~5.4(1)]{GJR26}, proving the forward implication.  Conversely,
suppose that \(\iota^*g\) is \(\mathbf R\)-PL.  By
\cite[Definition~5.1 and \S5.2(4)]{GJR26}, on a rational G-cover of
\(X_0\) it has presentations
\[
                  \iota^*g=\sum_{\ell=1}^r
                         a_\ell\log|\overline f_\ell|,
 \qquad a_\ell\in\mathbf R,\quad
        \overline f_\ell\in\mathcal O_{X_0}^{\times}.
\]
Lift the cover to \(X\) by the preceding paragraph.  Each
\(\overline f_\ell\) lifts to an analytic function \(f_\ell\).  It is a
unit: if \(\overline b_\ell\) is its inverse and \(b_\ell\) is a lift, then
\[
                         f_\ell b_\ell=1+n_\ell
\]
with \(n_\ell\) nilpotent, and
\(b_\ell(1+n_\ell)^{-1}\) is an inverse of \(f_\ell\).  The functions
\(f_\ell\) and \(\overline f_\ell\) have identical absolute values at the
identified analytic points.  Hence
\[
                         g=\sum_{\ell=1}^r
                              a_\ell\log|f_\ell|
\]
on the lifted cover.  Thus \(g\) is \(\mathbf R\)-PL.  This unit-lifting
calculation is also the nilpotent-immersion mechanism behind
\cite[Proposition~15.1.3(b)]{CLD25}.
\end{proof}

\begin{proof}[Proof of \Cref{lem:p061-product-factor-reduction}]
Choose strict presentations
\[
 A_i=L\langle\mathbf T_i\rangle/\mathfrak a_i,
 \qquad
 B=L\langle\mathbf U\rangle/\mathfrak b,
\]
and let \(\widetilde I_i\subset L\langle\mathbf T_i\rangle\) and
\(\widetilde I_B\subset L\langle\mathbf U\rangle\) be the inverse images
of \(I_i\) and \(I_B\).  In the joint Tate algebra
\[
 R=L\langle\mathbf T_1,\ldots,\mathbf T_m,\mathbf U\rangle
\]
put
\[
 \mathfrak h=\mathfrak a_1R+\cdots+\mathfrak a_mR+\mathfrak bR.
\]
The universal property of the strict affinoid completed tensor product
gives the explicit presentations
\[
 P\simeq R/\mathfrak h,
 \qquad
 Q\simeq R/\bigl(\mathfrak h+\widetilde I_1R+\cdots+
                      \widetilde I_mR+\widetilde I_BR\bigr).
\]
All displayed ideals are closed: Tate algebras and their affinoid quotients
are noetherian, and every ideal in an affinoid algebra is closed.  Thus the
map \(P\to Q\) is a strict surjection, and its kernel is exactly
\[
                         J=I_1P+\cdots+I_mP+I_BP.
\]
Choose \(e_i,e_B\geq1\) with \(I_i^{e_i}=0\) and \(I_B^{e_B}=0\).
Every monomial of degree
\[
                         E=1+\sum_{i=1}^m(e_i-1)+(e_B-1)
\]
in the \(m+1\) displayed summands of \(J\) contains at least \(e_i\)
factors from one \(I_iP\), or at least \(e_B\) factors from \(I_BP\).
Consequently \(J^E=0\).  This argument includes \(m=0\).

Every multiplicative seminorm on \(P\) kills \(J\), so the induced map
\(\mathcal M(Q)\to\mathcal M(P)\) is bijective.  It is a homeomorphism by
the compact-Hausdorff argument used in the preceding proof.  More
precisely, let \(V=\mathcal M(P_V)\) be a strict affinoid domain in
\(\mathcal M(P)\).  Pullback of the coherent closed immersion
\(\mathcal M(Q)\to\mathcal M(P)\) cuts out on \(V\) the coherent ideal
\(JP_V\); its inverse image is therefore
\[
                         \mathcal M(P_V/JP_V).
\]
This is the affinoid closed-subspace construction of
\cite[\S11.2.2]{CLD25}.  Since \((JP_V)^E=0\), the base-changed immersion
is again nilpotent and a homeomorphism.  This proves the final assertion
without assuming that an arbitrary affinoid domain is itself rational.
\end{proof}

\begin{proof}[Proof of
\Cref{thm:p061-arbitrary-curve-factors-ph-rpl}]
We first prove the assertion on the displayed product.  The question is
local.  Put \(Y_0=Y_{\mathrm{red}}\).  Quasi-smoothness makes local
dimension locally constant, so around the projection of the chosen point
we may shrink \(Y\), whose underlying topological space is that of \(Y_0\),
until \(Y_0\) is pure-dimensional.  Let \(K'\) be the completion of an
algebraic closure of \(K\), attach a prime to scalar extensions to \(K'\),
and put
\[
                              E=Y_0'.
\]
Quasi-smoothness is stable under complete scalar extension, so \(E\) is
quasi-smooth.  The function \(h'\) is classically pluriharmonic by
\cite[Proposition~7.2(6)]{GR25Classical}.

For every \(i\), put
\[
                         D_i=(C_i')_{\mathrm{red}}
\]
and set
\[
             W=D_1\times_{K'}\cdots\times_{K'}D_m\times_{K'}E.
\]
The natural map
\[
                              j:W\longrightarrow X'
\]
uses the actual reductions of the \(C_i'\) and the base change of the
nilpotent closed immersion \(Y_0\to Y\).  Thus no commutation of reduction
with scalar extension is assumed.  To check the product map locally, fix a
point of \(X'\), choose strict affinoid domains
\(\mathcal M(A_i)\subset C_i'\) around all projected curve points and a
strict affinoid domain \(\mathcal M(B)\subset Y'\) around the projected
\(Y'\)-point, and take their affinoid product.  By the closed-subspace and
analytic-domain compatibility recalled in \cite[\S11.2.2]{CLD25}, its
inverse image in \(D_i\) is
\(\mathcal M(A_i/N_i)\), where \(N_i\) is the nilradical of \(A_i\), and
its inverse image in \(E\) is \(\mathcal M(B/I_B)\), where \(I_B\) is the
restriction of the defining ideal of \(E\to Y'\).  Affinoid algebras are
noetherian, so all the \(N_i\) are nilpotent; the ideal \(I_B\) is
nilpotent because it is obtained by base change from the locally nilpotent
ideal defining \(Y_0\to Y\).  Hence
\Cref{lem:p061-product-factor-reduction} shows on these covering product
neighborhoods that \(j\) is a nilpotent closed immersion and a
homeomorphism, including for the rational G-topologies.  In particular,
\(W\) is good and strict.  We neither assert nor use that \(W\) is reduced,
that it equals \((X')_{\mathrm{red}}\), or that \(E=(Y')_{\mathrm{red}}\);
the nilpotence of \(j\) is the only required product/reduction statement.
The pullback
\[
                              h_W=j^*h'
\]
is classically pluriharmonic by the arbitrary-morphism functoriality
\cite[Proposition~7.2(7)]{GR25Classical}.

Let
\[
                         \nu_i:D_i^\nu\longrightarrow D_i
\]
be the normalization.  All arguments below are local, so we work on
affinoid product neighborhoods.  Analytic normalization is finite and is
compatible with restriction to analytic domains
\cite[\S11.3.1]{CLD25}.  The field \(K'\) is algebraically closed and hence
perfect, so \cite[Lemma~3.9(1)]{GR25Classical} makes every \(D_i^\nu\)
quasi-smooth.  For \(0\leq r\leq m\), put
\[
\begin{aligned}
 W_r={}&D_1\times_{K'}\cdots\times_{K'}D_r
       \times_{K'}D_{r+1}^\nu\times_{K'}\cdots\\
      &\hspace{37mm}\cdots\times_{K'}D_m^\nu\times_{K'}E.
\end{aligned}
\]
Let \(\rho_r:W_r\to W_m=W\) normalize the factors with index greater
than \(r\), and let \(h_r=\rho_r^*h_W\).  Finite products of the finite
normalization maps make every \(\rho_r\) finite.  Hence the spaces \(W_r\)
are good and strict locally because they are finite over the good strict
space \(W\).  Every \(h_r\) is classically pluriharmonic by
\cite[Proposition~7.2(7)]{GR25Classical}.  The space \(W_0\), being a
finite product of quasi-smooth spaces, is quasi-smooth.  Hence the
production-verified
\Cref{thm:p061-quasismooth-ph-rpl} gives
\[
                         h_0\in\operatorname{PL}_{\mathbf R}(W_0).
\]
When \(m=0\), this says simply that \(W_0=W=E\) and \(h_0=h_W\); the
descent induction below is empty.

For \(1\leq r\leq m\), let \(q_r:W_{r-1}\to W_r\) be induced by
\(\nu_r\) on the \(r\)-th factor.  Then
\[
                              q_r^*h_r=h_{r-1}.
\]
Fix a point of \(W_r\).  In all factors other than \(D_r\), restrict to
strict affinoid neighborhoods on which the parameter product is
pure-dimensional.  Such neighborhoods exist: analytic curves are separated
and pure of dimension one by the curve convention, their normalizations
have the same dimension, \(E\) was made locally pure-dimensional, and the
strict affinoid product dimension formula makes a finite product of these
neighborhoods pure of the sum of their dimensions.  Dimension invariance
under reduction and complete scalar extension, as well as componentwise
purity, are recalled in \cite[\S11.2.3--\S11.2.5]{CLD25}.  Normalization
restricts to the inverse image of the chosen curve domain.  The hypotheses
of the already production-verified
\Cref{prop:p061-curve-product-normalization-descent} therefore hold on this
local product.  Applied over the algebraically closed field \(K'\), it
descends \(\mathbf R\)-piecewise linearity from \(h_{r-1}\) to \(h_r\).
Since piecewise linearity is local, doing this at every point and inducting
on \(r\) gives
\[
                         h_W=h_m\in
                         \operatorname{PL}_{\mathbf R}(W).
\]
This use of
\Cref{prop:p061-curve-product-normalization-descent} is not circular: that
proposition was proved independently for a reduced curve over an
algebraically closed field, using
\Cref{lem:p061-rational-section-tube,%
lem:p061-parameterized-rigid-point}, before any geometric-reducedness
theorem.

Apply \Cref{lem:p061-nilpotent-rpl-invariance} to \(j:W\to X'\).
It yields \(h'\in\operatorname{PL}_{\mathbf R}(X')\).  The extension
\(K'/K\) is the completion of an algebraic extension, so ground-field
descent \cite[Corollary~5.12]{GJR26} gives
\[
                         h\in\operatorname{PL}_{\mathbf R}(X).
\]
Finally, both \(h\) and \(-h\) are classically psh.  For an
\(\mathbf R\)-PL function on a good strict space,
\cite[Theorem~4.3]{GR25Classical} identifies classical psh with
semipositivity.  Both signs are therefore semipositive, which is
harmonicity by \cite[Definition~7.2]{GJR26}.

It remains to justify the open-subspace assertion.  It is enough to treat
an open analytic subspace \(U\) of the displayed product.  Work after
extension to \(K'\) and pull back the nilpotent immersion \(W\to X'\) to
\[
                         U_W=U'\times_{X'}W\longrightarrow U'.
\]
This is again a nilpotent closed immersion and a homeomorphism.  Fix
\(u\in U_W\), and let \(I\) be the set of curve indices for which the
projection of \(u\) to \(D_i\) is nonnormal.  If \(I\) is empty, a
neighborhood of \(u\) is quasi-smooth and
\Cref{thm:p061-quasismooth-ph-rpl} applies.

If \(I\ne\varnothing\), every corresponding nonnormal point is rigid by
\cite[Lemma~3.9(2)]{GR25Classical}, and hence \(K'\)-rational.  First
shrink every remaining curve coordinate inside its normal locus.  Apply
the neighborhood-basis part of
\Cref{lem:p061-rational-section-tube} successively to the rational
coordinates indexed by \(I\).  It supplies strict affinoid neighborhoods
\(U_i\) of those curve points and a strict affinoid neighborhood \(V\) in
the product of the remaining coordinates such that
\[
                         \prod_{i\in I}U_i\times_{K'}V
                         \subset U_W.
\]
Normal reduced curves over the perfect field \(K'\) are quasi-smooth by
\cite[Lemma~3.9(1)]{GR25Classical}, so \(V\) is quasi-smooth.  The product
assertion already proved, applied to the curve domains \(U_i\) and the
quasi-smooth factor \(V\), proves \(\mathbf R\)-piecewise linearity near
\(u\).  Thus the pullback of the original function is \(\mathbf R\)-PL
on \(U_W\).  Nilpotent invariance gives the result on \(U'\), and
\cite[Corollary~5.12]{GJR26} descends it to \(U\).
Both signs of the original function on \(U\) remain classically psh.
Now \cite[Theorem~4.3]{GR25Classical} makes both signs semipositive, so the
function is Gubler--Jell--Rabinoff harmonic by
\cite[Definition~7.2]{GJR26}.  Harmonicity and piecewise linearity are
local, which also proves the formulation for every good strict space
analytically locally isomorphic to such open subspaces.

This proves exactly the scope stated in the theorem.  No step descends
through a general higher-dimensional nonflat normalization: finite
nonflat descent is used only in the verified curve-normalization
proposition, where the rational-section tube supplies the missing analytic
retraction.
\end{proof}
